\exercisesection{§ 4 的习题}

若 \(\mathbb{R}\) 是根系，以 \(W^{+}(R)\) 表示 \(W(R)\) 中行列式为 1 的元素所成的集合。

¶ 1) 令 R 为 E8 型根系。

\begin{enumerate}
\def\labelenumi{\alph{enumi})}
\item
  证明，若 \(\alpha, \beta \in \mathbb{R}\) 模 \(2\mathrm{Q}(\mathbb{R})\) 同余，则 \(\beta = \pm \alpha\)。
\item
  由 a) 推出，若 \(w \in \mathrm{W}(\mathrm{R})\) 在 \(\mathrm{Q}(\mathrm{R}) / 2\mathrm{Q}(\mathrm{R})\) 上平凡作用，则 \(w = \pm 1\)。
\item
  用第 10 号的记号，证明 Q(R) 上的二次型 \(\frac{1}{2}(x|x)\) 通过取商定义 \(q_{8}\) 于 \(\mathbf{F}_2\)-向量空间 Q(R)/2Q(R) 上。证明 \(q_{8}\) 的伪判别式（参见《代数》第 IX 章 §9，习题 9）为零，且 \(q_{8}\) 的指标为 4。
\item
  令 \(\mathbf{O}(q_{8})\) 为关于此型的 \(\mathrm{Q}(\mathrm{R})/2\mathrm{Q}(\mathrm{R})\) 的正交群。通过取商定义一个同态
\end{enumerate}

\[
h: \mathrm{W} (\mathrm{R}) \rightarrow \mathbf {O} (q _ {8}).
\]

证明（通过比较阶）序列

\[
1 \rightarrow \{1, - 1 \} \rightarrow \mathrm{W(R)} \stackrel {h} {\longrightarrow} \mathbf {O} (q _ {8}) \rightarrow 1
\]

是正合的。

\begin{enumerate}
\def\labelenumi{\alph{enumi})}
\setcounter{enumi}{4}
\tightlist
\item
  证明 \(\mathrm{W}^{+}(\mathrm{R})\) 在 \(h\) 下的像是《代数》第 IX 章 §9 第 5 号中定义的 \(\mathbf{O}^{+}(q_{8})\) 的子群（它属于 \(\mathbf{O}(q_8)\)）。推出 \(\mathrm{W}^{+}(\mathrm{R}) / \{1, -1\}\) 是一个非阿贝尔单群。
\end{enumerate}

\begin{enumerate}
\def\labelenumi{\arabic{enumi})}
\setcounter{enumi}{1}
\tightlist
\item
  令 R 为 \(E_{6}\) 型根系。
\end{enumerate}

\begin{enumerate}
\def\labelenumi{\alph{enumi})}
\item
  置 \(E = Q(R)/3P(R)\)。这是一个 5 维 \(F_{3}\)-向量空间。用第 12 号的记号，证明标量积 \((x|y)\) 在 E 上定义一个非退化对称双线性型 \(\varphi\)。证明 R 的任意两个不同元素在 E 中的像不同。
\item
  令 \(\mathbf{O}(\varphi)\) 为 \(\varphi\) 的正交群。则 \(\mathbf{O}(\varphi) = \{1, -1\} \times \mathbf{SO}(\varphi)\)。自旋范数定义了从 \(\mathbf{SO}(\varphi)\) 到 \(\{1, -1\}\) 的满同态，其核记为 \(\mathbf{SO}^{+}(\varphi)\)。群 \(\mathbf{SO}^{+}(\varphi)\) 是阶为 25920 的单群。商 \(\mathbf{O}(\varphi)/\mathbf{SO}^{+}(\varphi)\) 的类型为 (2, 2)。推出 \(\mathbf{O}(\varphi)\) 含有唯一的指数为 2 的子群 \(\Omega(\varphi)\)，它不同于 \(\mathbf{SO}(\varphi)\) 且不含 \(-1\)。
\item
  A(R) 的每个元素通过取商定义 \(\mathbf{O}(\varphi)\) 的一个元素。证明（通过比较阶）这给出从 A(R) 到 \(\mathbf{O}(\varphi)\) 的同态。W(R) 在此同态下的像为 \(\Omega(\varphi)\)，而 \(\mathrm{W}^{+}(\mathbb{R})\) 的像为 \(\mathbf{SO}^{+}(\varphi)\)。因此，\(\mathrm{W}(\mathbb{R})\) 是一个以阶为 25920 的单群为被扩张对象、商为 \(\mathbf{Z}/2\mathbf{Z}\) 的扩张。
\item
  令 \(F = Q(R)/2Q(R)\)。这是一个 6 维 \(F_{2}\)-向量空间。证明二次型 \(\frac{1}{2}(x|x)\) 通过取商在 F 上定义一个伪判别式等于 1 的非退化二次型 \(q_{6}\)。若以 \(\mathbf{O}(q_{6})\) 表示相应的正交群，则通过取商定义一个同态
\end{enumerate}

\[
h: \mathrm{W} (\mathrm{R}) \rightarrow \mathbf {O} (q _ {6}).
\]

证明 \(h\) 是单射（注意 \(-1 \notin \mathrm{W}(\mathbb{R})\)），再证明它是同构（比较阶）。推出存在从 \(\mathrm{W}^{+}(\mathbb{R})\) 到 \(\mathbf{O}^{+}(q_{6})\) 的同构（参见《代数》第 IX 章 §9 第 5 号）。

\begin{enumerate}
\def\labelenumi{\alph{enumi})}
\setcounter{enumi}{2}
\tightlist
\item
  通过比较 c) 和 d)，证明 \(^{4}\) \(\mathbf{SO}^{+}(\varphi)\) 同构于 \(\mathbf{O}^{+}(q_{6})\)。
\end{enumerate}

¶ 3) 令 R 为 \(E_{7}\) 型根系。

\begin{enumerate}
\def\labelenumi{\alph{enumi})}
\item
  置 \(E = Q(R)/2P(R)\)。这是一个 6 维 \(F_{2}\)-向量空间。用第 11 号的记号，证明标量积 \((x|y)\) 在 E 上定义一个非退化交错双线性型。
\item
  由 a) 推出存在正合序列
\end{enumerate}

\[
1 \rightarrow \{1, - 1 \} \rightarrow \mathrm{W} (\mathrm{R}) \xrightarrow {h} \mathbf {S p} (6, \mathbf {F} _ {2}) \rightarrow 1.
\]

（利用 \(\mathbf{Sp}(6,\mathbf{F}_2)\) 的阶为 \(2^{9}.3^{4}.5.7\) 这一事实。）

\begin{enumerate}
\def\labelenumi{\alph{enumi})}
\setcounter{enumi}{2}
\tightlist
\item
  证明 \(h\) 在 \(\mathrm{W}^{+}(\mathrm{R})\) 上的限制是从 \(\mathrm{W}^{+}(\mathrm{R})\) 到 \(\mathbf{Sp}(6, \mathbf{F}_2)\) 的同构。
\end{enumerate}

\(d\) ) 给出 \(b\) ) 的第二个证明，利用二次型 \(q_{7}\)；该二次型定义在 \(\mathbf{F}_2\)-向量空间 \(\mathrm{Q(R) / 2Q(R)}\) 上，由 \(\frac{1}{2} (x|x)\) 诱导，并利用同构

\[
\mathbf {O} (q _ {7}) \rightarrow \mathbf {S p} (6, \mathbf {F} _ {2}).
\]

\begin{enumerate}
\def\labelenumi{\arabic{enumi})}
\setcounter{enumi}{3}
\tightlist
\item
  令 R 为 V 中的不可约约化根系。\((\alpha_{1},\ldots,\alpha_{l})\) 为 R 的一个基。\(\tilde{\alpha}\) 为最高根。置 \(\tilde{\alpha}=n_{1}\alpha_{1}+\cdots+n_{l}\alpha_{l}\)。我们将确定 R 中不同于 R 的极大闭对称子集。
\end{enumerate}

\begin{enumerate}
\def\labelenumi{\alph{enumi})}
\item
  令 \(i \in \{1, 2, \ldots, l\}\)。令 \(R_i\) 为满足 \(\alpha \in R\) 且由 \(\alpha_j\)（\(j \neq i\)）作线性组合的元素所成的集合。证明 \(R_i\) 极大，当且仅当 \(n_i = 1\)。
\item
  令 \(i \in \{1, 2, \ldots, l\}\)，并假设 \(n_i > 1\)。令 \(S_i\) 为根 \(\sum_{j=1}^{l} m_j \alpha_j\) 中满足 \(m_i \equiv 0 \pmod{n_i}\) 的根所成的集合。证明 \(S_i\) 极大，当且仅当 \(n_i\) 为素数。（若 \(n_i = ab\)，其中 a \textgreater{} 1、b \textgreater{} 1，考虑 R 的子集 \(S'\)，其元素为根 \(\sum_{j} m_j \alpha_j\)，并满足 \(m_i \equiv 0 \pmod{a}\)，并证明 \(S'\) 真包含 \(S_i\)。）证明根 \(-\tilde{\alpha}, \alpha_j (j \neq i)\) 构成 \(S_i\) 的一个基（其秩为 l）。推出 \(S_i\) 的 Dynkin 图。
\item
  R 的每个极大闭对称子集都可由 W(R) 的一个元素变换为 a) 或 b) 中描述的某个子集。（令 \(\Sigma\) 为这样的子集。则 \(\Sigma = R \cap H\)，其中 H 是 Q(R) 的有限指数子群（§1，习题 6 b)）；可设 Q(R)/H 是循环的。于是存在 \(u^{*} \in V^{*}\)，使得 \(\Sigma\) 是由满足 \(\alpha \in R\) 且 \(\langle u^{*}, \alpha \rangle \in Z\) 的元素组成的集合。）
\item
  列出不同类型的不可约约化根系 R 的极大闭对称子集。
\end{enumerate}

\begin{enumerate}
\def\labelenumi{\arabic{enumi})}
\setcounter{enumi}{4}
\tightlist
\item
  令 R 为根系，令 \(P'(R)\) 为 §1 习题 8 中引入的 \(P(R)\) 子群。证明 \(P'(R) = P(R)\) 当 R 为 \(A_l\) 型且 l 为偶数，或为 \(B_l\) 型且 \(l \equiv 0,3 \pmod{4}\)，或为 \(D_l\) 型且 \(l \equiv 0,1 \pmod{4}\)，或为 \(G_2\) 型、\(F_4\) 型、\(E_6\) 型或 \(E_8\) 型时成立。证明 \(P'(R) = Q(R)\) 当 R 为 \(C_l\) 型，或为 \(B_l\) 型且 \(l \equiv 1,2 \pmod{4}\)，或为 \(E_7\) 型，或为 \(A_l\) 型时成立。若 R 为 \(A_l\) 型且 l 为奇数且 \textgreater1，则 \(P'(R)/Q(R)\) 是循环群 \(P(R)/Q(R)\) 中唯一的指数为 2 的子群。若 R 为 \(D_l\) 型且 \(l \equiv 2,3 \pmod{4}\)，则 \(P'(R)/Q(R)\) 是 \(P(R)/Q(R)\) 中在 \(\Lambda(R)\) 下稳定的唯一的阶为 2 的子群。
\end{enumerate}

¶ 6) 令 R 为不可约约化根系，\((\alpha_{1},\ldots,\alpha_{l})\) 为 R 的一个基，\(p_{1}\alpha_{1}+\cdots+p_{l}\alpha_{l}\) 为最高根，\(a_{1}\alpha_{1}+\cdots+a_{l}\alpha_{l}\) 为正根之和，\(m_{1},\ldots,m_{l}\) 为 W(R) 的指数，g 为其阶，f 为连通指标。

\begin{enumerate}
\def\labelenumi{\alph{enumi})}
\tightlist
\item
  证明在每种情形下，
\end{enumerate}

\[
l! p _ {1} p _ {2} \dots p _ {l} m _ {1} m _ {2} \dots m _ {l} = a _ {1} a _ {2} \dots a _ {l}.
\]

\begin{enumerate}
\def\labelenumi{\alph{enumi})}
\setcounter{enumi}{1}
\tightlist
\item
  证明
\end{enumerate}

\[
g m _ {1} m _ {2} \dots m _ {l} = f a _ {1} a _ {2} \dots a _ {l}.
\]

（利用 \(a\) ）以及 §2 第 4 号的命题 7。）

\begin{enumerate}
\def\labelenumi{\alph{enumi})}
\setcounter{enumi}{2}
\tightlist
\item
  对任意正根 \(\alpha = \sum_{i} c_{i} q_{i}\)，令 \(c(\alpha) = \sum_{i} c_{i}\)。计算每种情形下的多项式
\end{enumerate}

\[
\mathrm{P} (t) = \sum_ {\alpha > 0} t ^ {c (\alpha)}.
\]

证明 \(^{5}\) \(P(t) = t \sum_{i=1}^{l} (1 + t + \cdots + t^{m_{i}-1})\)。

\begin{enumerate}
\def\labelenumi{\arabic{enumi})}
\setcounter{enumi}{6}
\tightlist
\item
  令 R 为不可约约化根系。
\end{enumerate}

\begin{enumerate}
\def\labelenumi{\alph{enumi})}
\item
  证明从 \(A(R)/W(R)\) 到 \(P(R)/Q(R)\) 的自同构群的典范同态是单射。
\item
  推出 -1 属于 W(R)，当且仅当 Q(R) ⊆ 2P(R)。
\end{enumerate}

\begin{enumerate}
\def\labelenumi{\arabic{enumi})}
\setcounter{enumi}{7}
\item
  令 \(R_{1}\) 和 \(R_{2}\) 为两个不可约约化根系。证明若 \(W(R_{1})\) 和 \(W(R_{2})\) 的阶相同，则 \(R_{1}\) 同构于 \(R_{2}\) 或 \(R_{2}^{-}\)。（利用分类。）若不假设 \(R_{1}\) 不可约，该结果仍然成立吗？
\item
  令 R 为秩 l 的根系，p 为整除 A(R) 的阶的素数。证明 \(p \leqslant l + 1\)。（化归到不可约情形，并利用分类。）
\item
  \begin{enumerate}
  \def\labelenumii{\alph{enumii})}
  \tightlist
  \item
    令 \((W, S)\) 为不可约有限 Coxeter 系统。置
  \end{enumerate}
\end{enumerate}

\[
W (t) = \sum_ {w \in W} t ^ {l (w)} \quad (\text { cf.   Chap.   IV,   } \S 1, \text {   Exerc.   26. })
\]

令 \(m_{1},\ldots ,m_{l}\) 为 W 的指数（第 V 章 §6 第 2 号）。证明公式

\[
\mathrm{W} (t) = \prod_ {i = 1} ^ {l} \left(1 + t + \dots + t ^ {m _ {i}}\right)
\]

对 \(l\) 的较小值成立（利用第 IV 章 §1 的定理 1 和习题 26）\(^{6}\)。

\begin{enumerate}
\def\labelenumi{\alph{enumi})}
\setcounter{enumi}{1}
\tightlist
\item
  令 R 为约化不可约根系，\(W\)（分别为 \(W_{a}\)）为 R 的 Weyl 群（分别为仿射 Weyl 群），并赋予由所选胞腔决定的 Coxeter 群结构。按上文定义 \(W(t)\) 和指数 \(m_{i}\)，并置
\end{enumerate}

\[
W _ {a} (t) = \sum_ {w \in W _ {a}} t ^ {l (w)}.
\]

证明公式

\[
\mathrm{W} _ {a} (t) = \mathrm{W} (t) \prod_ {i = 1} ^ {l} \frac {1}{1 - t ^ {m _ {i}}} = \prod_ {i = 1} ^ {l} \frac {1 + t + \cdots + t ^ {m _ {i}}}{1 - t ^ {m _ {i}}}
\]

对 \(l\) 的较小值成立（利用第 IV 章 §1 的定理 2 和习题 26）\(^{7}\)。

¶ 11) 令 (W.S) 为 \(H_{3}\) 型 Coxeter 系统（参见定理 1）。

\begin{enumerate}
\def\labelenumi{\alph{enumi})}
\item
  采用第 V 章 §6 第 2 号引理 2 的证明中的记号，证明 \((z', z'') = \pi/5\)；推出 \(h\)（群 \(W\) 的 Coxeter 数）等于 10，且 \(W\) 的指数为 1、5 和 9。
\item
  利用 a) 证明 \(\operatorname{Card}(W) = 120\)，并证明 W 中反射的数目为 15。
\item
  通过应用第 V 章 §3 的习题 5，重新得到公式 \(\operatorname{Card}(W) = 120\)。d) 令 \(\mathcal{A}_5\) 为 \(\{1, \ldots, 5\}\) 的交错群；若 \(a, b, c, d\) 是 \(\{1, \ldots, 5\}\) 的互不相同的元素，以 \((ab)\) 表示交换 \(a\) 与 \(b\) 的换位，以 \((ab)(cd)\) 表示换位 \((ab)\) 与 \((cd)\) 的乘积。令
\end{enumerate}

\[
r _ {1} = (14) (23), \quad r _ {2} = (12) (45), \quad r _ {3} = (12) (34).
\]

证明 \((r_1r_2)^5 = (r_2r_3)^3 = (r_1r_3)^2 = 1\)。推出存在一个同态 \(f: \mathrm{W} \to \mathcal{A}_5\)，将 S 映到 \(\{r_1, r_2, r_3\}\)；证明 \(f\) 是满射。\\
e) 令 \(\varepsilon: \mathrm{W} \to \{\pm 1\}\) 为同态 \(w \mapsto (-1)^{l(w)}\)。证明

\[
(f, \varepsilon): W \rightarrow \mathcal {A} _ {5} \times \{\pm 1 \}
\]

是一个同构。（利用所考虑的两个群阶相同这一事实。）

¶ 12) 令 \((1, i, j, k)\) 为四元数域 H 的典范基，借此将 H 识别为 \(\mathbf{R}^4\)。赋予 H 标量积 \(\frac{1}{2}(x\bar{y} + y\bar{x})\)。令 \(\Gamma\) 为范数为 1 的四元数乘法群。

\begin{enumerate}
\def\labelenumi{\alph{enumi})}
\item
  若 \(a \in \Gamma\)，则 H 中将 a 变为 -a 的正交反射 \(s_{a}\) 是映射 \(x \mapsto -a\bar{x}a\)。
\item
  令 \(q = \cos \frac{\pi}{5} - \frac{1}{2}i + (\cos \frac{3\pi}{5})j \in \Gamma\)，且 \(r = \frac{1}{2}(1 + i + j + k) \in \Gamma\)。令 Q 为从 1、q、r 对坐标作任意偶置换和符号改变所得到的四元数集合。则 Q 是 \(\Gamma\) 的阶为 120 的子群。
\item
  令 W 为 \(\mathbf{GL}(\mathbf{H})\) 中由 \(s_{a}\)（\(a \in Q\)）生成的子群。证明 W 保持 Q 稳定，且是有限的、不可约的和非结晶的；推出 W 为 \(H_{4}\) 型（参见定理 1）。
\item
  证明 W 在 Q 上传递作用（利用第 IV 章 §1 的命题 3）。
\item
  令 \(a_0 \in Q\)，令 \(V\) 为 \(\mathbf{H}\) 中与 \(a_0\) 正交的向量子空间，令 \(W_0\) 为稳定 \(a_0\) 的 \(W\) 的子群。证明 \(W_0\) 在 \(V\) 上的限制是由反射生成的不可约群（第 V 章 §3，第 3 号），且它是非结晶的。推出 \(W_0\) 是 \(H_3\) 型 Coxeter 群。
\item
  证明 \(\operatorname{Card}(W)=2^{6}3^{2}5^{2}\)（利用 (d)、(e) 和习题 11）。
\end{enumerate}

g) 证明 \(\mathbf{W}\) 的 Coxeter 数等于 30，且其指数为 1、11、19、29。

\begin{enumerate}
\def\labelenumi{\arabic{enumi})}
\setcounter{enumi}{12}
\tightlist
\item
  令 V 为 \(R^{9}\) 中方程 \(x_{1} + \cdots + x_{9} = 0\) 所确定的超平面。令 R 为 V 中由下列点组成的子集
\end{enumerate}

\[
(2. 2, 2. - 1, \dots 1. - 1, \dots 1, \dots 1, \dots 1), \quad (- 2. - 2. - 2. 1, 1. 1. 1. 1. 1),
\]

\[
(3. - 3. 0. 0. 0. 0. 0. 0. 0)
\]

以及通过置换坐标从这些点得到的点。证明 R 是 V 中的 \(E_{8}\) 型根系。

\begin{enumerate}
\def\labelenumi{\arabic{enumi})}
\setcounter{enumi}{13}
\item
  用第 7 号的记号，证明 \(\mathbf{R}^{l+1}\) 上将 \(\varepsilon_1\) 变为 \(\varepsilon_2\)、将 \(\varepsilon_2\) 变为 \(\varepsilon_3\)、\(\ldots\)、将 \(\varepsilon_l\) 变为 \(\varepsilon_{l+1}\) 并将 \(\varepsilon_{l+1}\) 变为 \(\varepsilon_1\) 的自同构，诱导出 R（\(A_l\) 型系统）的 Coxeter 变换。
\item
  确定每种约化不可约根系的极小权（§ 1，习题 24）。（可得到基本权 \(\omega_{1},\ldots,\omega_{l}\) 作为 \(A_{l}\) 型的权，\(\omega_{l}\) 作为 \(B_{l}\) 型的权，\(\omega_{1}\) 作为 \(C_{l}\) 型的权，\(\omega_{1},\omega_{l-1},\omega_{l}\) 作为 \(D_{l}\) 型的权，\(\omega_{1}\) 和 \(\omega_{6}\) 作为 \(E_{6}\) 型的权，\(\omega_{7}\) 作为 \(E_{7}\) 型的权，而 \(E_{8},F_{4}\) 和 \(G_{2}\) 型没有极小权。）
\end{enumerate}

¶ 16) 令 (W.S) 为不可约有限 Coxeter 系统，且 n = Card(S)。

\begin{enumerate}
\def\labelenumi{\alph{enumi})}
\item
  若 \((W, S)\) 不是 \(F_4\) 型，证明存在 \(X\)，它是 \(S\) 的含有 \(n - 1\) 个元素的子集，使得 \((W_X, X)\) 是 \(A_{n-1}\) 型。
\item
  通过典范表示将 W 识别为 \(\mathbf{GL}(\mathbf{R}^{S})\) 的一个子群（第 V 章 §4）。证明存在一个基 \((c_{1},\ldots,c_{n})\)，它是 \(R^{S}\) 的基，使得对每个置换 \(\sigma\in S_{n}\)，在 \(R^{S}\) 上将 \(e_{i}\) 变为 \(c_{\sigma(i)}\)（\(1\leqslant i\leqslant n\)）的自同构属于 W（``Burnside 定理''）。（当 (W.S) 不是 \(F_{4}\) 型时，利用 a)；当它是 \(F_{4}\) 型时，注意 W 含有一个 \(D_{4}\) 型子群（参见第 9 号），从而将问题化归为前一种情形。）
\item
  令 E 为 (W.S) 的自同构群的一个子群。证明 E 由 W 半直积得到的 E.W 以典范方式嵌入 \(\mathbf{GL}(\mathbf{R}^{\mathrm{S}})\)。证明如此定义的 \(\mathbf{GL}(\mathbf{R}^{\mathrm{S}})\) 的子群由反射生成，但以下四种情形除外：
\end{enumerate}

\begin{enumerate}
\def\labelenumi{(\roman{enumi})}
\item
  (W.S) 是 \(A_{n}\) 型，\(n \geqslant 4\)：群 \(E\) 的阶为 2。
\item
  (W.S) 是 \(\mathrm{D}_4\) 型；群 \(E\) 的阶为 3。
\item
  (W.S) 是 \(F_4\) 型；群 \(E\) 的阶为 2。
\item
  (W.S) 是 \(\mathbf{E}_6\) 型；群 \(E\) 的阶为 2。
\end{enumerate}

证明在情形 (i) 和 (iv) 中，E.W = \(\{\pm1\} \times W\)。证明在情形 (iii) 中，群 E.W 不保持 \(R^{S}\) 中的任何格稳定。

\begin{enumerate}
\def\labelenumi{\alph{enumi})}
\setcounter{enumi}{3}
\tightlist
\item
  令 \(W_{1}\) 为由反射生成的 \(\mathbf{GL}_{n}(\mathbf{R})\) 的有限子群，并假设 \(W_{1}\) 不可约且本质。令 G 为 \(\mathbf{GL}_{n}(\mathbf{R})\) 中包含 \(W_{1}\) 的有限子群。证明 G 或者由反射生成，或者具有 E.W 的形式，其中 E 和 W 属于 c) 中的类型 (i)、(ii)、(iii)、(iv) 之一。
\end{enumerate}

（令 W 为属于 G 的反射所生成的群。群 G 置换相对于 W 的各个室。像 §2 第 3 号中那样推出 G 具有上述 E.W 的形式。）
% semantic-fragments: legacy-input-seam
\relax{}
